\begin{lemma}
Assume \(\noderef{ParameterHierarchy}(\varepsilon,\varepsilon_1,\varepsilon_2,n_0)\),
\(\eta>0\), and
\[
\eta\le
\frac{\varepsilon(1-\varepsilon-22\varepsilon^2)}{320(1+\varepsilon)}.
\]
For all sufficiently large \(n\), with
\[
K=\noderef{TwoBiteNaturalIndependenceScale}(1+\varepsilon,n),
\qquad
p=\noderef{TwoBiteEdgeProbability}\!\left(\frac12,n\right),
\]
the following holds uniformly for \(0\le\ell_R,\ell_B\le K\).  Let
\[
x_R=\ell_R/K,\qquad x_B=\ell_B/K.
\]
If
\[
1-\varepsilon/2<x_R+x_B<1+\varepsilon/2,\qquad x_B\le x_R,
\]
then \(pn<K-\ell_B\), and
\[
\begin{aligned}
&\left(\eta+\left(\frac12+\eta\right)-\frac{2-x_R-x_B}{2}\right)K\log n\\
&\quad
-p\left(
\binom{\ell_R}{2}_{\mathbb R}+\binom{\ell_B}{2}_{\mathbb R}
-\binom{K-\ell_R}{2}_{\mathbb R}
-\binom{pn}{2}_{\mathbb R}
-\binom{K-\ell_B-pn}{2}_{\mathbb R}
-\varepsilon^3K^2\right)\\
&\le
\frac{\varepsilon(-1+\varepsilon+22\varepsilon^2)}{16(1+\varepsilon)}
K\log n+3\eta K\log n.
\end{aligned}
\]
\end{lemma}
\begin{proof}
Since \(x_B\le x_R\) and \(x_R+x_B<1+\varepsilon/2\),
\[
x_B<\frac12+\frac{\varepsilon}{4}. \tag{1}
\]
For the rest of the proof write
\[
L=\log n,\qquad S=\sqrt{nL}.
\]
After increasing the threshold, \(L>0\) and \(S>0\).  The rounded comparison in
\(\noderef{ParameterHierarchyEventualComparisons}\) gives a number
\(\theta\) with
\[
0\le\theta<1,\qquad K=(1+\varepsilon)S+\theta. \tag{2}
\]
The definition
\[
p=\frac12\sqrt{\frac{\log n}{n}}
\]
then gives
\[
p=\frac{L}{2S},\qquad pn=\frac S2. \tag{3}
\]
Hence, with \(q=pn/K\),
\[
q=\frac{S}{2K}\le \frac{1}{2(1+\varepsilon)}.
\]
Since
\[
\frac12-\frac{\varepsilon}{4}-\frac{1}{2(1+\varepsilon)}
=\frac{\varepsilon(1-\varepsilon)}{4(1+\varepsilon)}>0,
\]
we have \(q<1/2-\varepsilon/4\).  From (1),
\(1-x_B>1/2-\varepsilon/4\), so \(q<1-x_B\).  Multiplying by the positive
quantity \(K\) gives
\[
K-\ell_B=(1-x_B)K>pn.
\]

Now expand the real binomials.  Since \(\ell_R=x_RK\),
\[
\begin{aligned}
\binom{\ell_R}{2}_{\mathbb R}
-\binom{K-\ell_R}{2}_{\mathbb R}
&=\frac{x_RK(x_RK-1)}2
-\frac{(1-x_R)K((1-x_R)K-1)}2\\
&=\frac{x_R^2K^2-x_RK-(1-x_R)^2K^2+(1-x_R)K}{2}\\
&=\frac{(2x_R-1)(K^2-K)}{2}. \tag{4}
\end{aligned}
\]
For the blue expression, write \(q=pn/K\).  Then
\[
\begin{aligned}
&\binom{\ell_B}{2}_{\mathbb R}
-\binom{pn}{2}_{\mathbb R}
-\binom{K-\ell_B-pn}{2}_{\mathbb R}\\
&=
\frac{K^2}{2}\left(x_B^2-q^2-(1-x_B-q)^2\right)
+\frac{K}{2}\left(-x_B+q+1-x_B-q\right)\\
&=
\frac{K^2}{2}\left(2x_B-1+2q-2x_Bq-2q^2\right)
+\frac{K}{2}(1-2x_B). \tag{5}
\end{aligned}
\]
Adding (4) and (5), subtracting \(\varepsilon^3K^2\), and writing
\[
t=x_R+x_B-1,\qquad q=\frac{pn}{K},
\]
gives the exact identity
\[
\begin{aligned}
&\binom{\ell_R}{2}_{\mathbb R}+\binom{\ell_B}{2}_{\mathbb R}
-\binom{K-\ell_R}{2}_{\mathbb R}
-\binom{pn}{2}_{\mathbb R}
-\binom{K-\ell_B-pn}{2}_{\mathbb R}
-\varepsilon^3K^2\\
&\qquad =
K^2\bigl(t+q(1-x_B)-q^2-\varepsilon^3\bigr)
+K(1-x_R-x_B). \tag{6}
\end{aligned}
\]
Put
\[
\alpha=\frac{1+\varepsilon}{2},
\qquad
q_0=\frac{1}{2(1+\varepsilon)},
\qquad
r=\frac{pK}{L}.
\]
Equations (2) and (3) give the two exact rounding formulae
\[
r=\frac{pK}{L}
=\frac{L}{2S}\cdot\frac{(1+\varepsilon)S+\theta}{L}
=\alpha+\frac{\theta}{2S},
\]
and
\[
0\le q_0-q
=\frac{1}{2(1+\varepsilon)}
-\frac{S}{2((1+\varepsilon)S+\theta)}
=\frac{\theta}{2(1+\varepsilon)K}.
\]
Since \(S\to\infty\), \(K\to\infty\), and \(L>0\), we may increase \(N\) so
that
\[
\frac{1}{2S}\le\frac{\eta}{16},\qquad
\frac{1}{2(1+\varepsilon)K}\le\frac{\eta}{16},\qquad
\frac{p}{L}=\frac{1}{2S}\le\frac{\eta}{16}. \tag{7}
\]
The estimates in (7), together with \(0\le\theta<1\), imply uniformly in the
present range
\[
|r-\alpha|\le\frac{\eta}{16},\qquad
|q-q_0|\le\frac{\eta}{16},\qquad
0\le q\le1,\qquad
p\le\frac{\eta}{16}\log n,
\tag{8}
\]
and \(K\log n\ge0\).  Let
\[
D=t+q(1-x_B)-q^2-\varepsilon^3,\qquad
D_0=t+q_0(1-x_B)-q_0^2-\varepsilon^3.
\]
Because \(|t|\le1\), \(0\le q\le1\), \(0\le x_B\le1\), and
\(0<\varepsilon<1\), we have \(|D|\le4\).  Also,
\[
|D-D_0|
\le |q-q_0|(1-x_B)+|q^2-q_0^2|
\le 3|q-q_0|,
\]
because \(0\le x_B,q,q_0\le1\).  Using (6), the exponent on the
left side of the lemma is
\[
\left(2\eta+\frac{t}{2}-rD\right)K\log n
-pK(1-x_R-x_B).
\]
Since
\[
-rD\le -\alpha D_0
+|r-\alpha|\,|D|
+\alpha |D-D_0|,
\]
and since \(\alpha<1\) and \(|1-x_R-x_B|\le1\), (8) gives the explicit
uniform bound
\[
\left(C_0+3\eta\right)K\log n, \tag{9}
\]
where
\[
C_0=\frac{t}{2}-\alpha\bigl(t+q_0(1-x_B)-q_0^2-\varepsilon^3\bigr).
\]
Indeed, the non-limiting errors in (9) are bounded by
\[
2\eta+\frac{4\eta}{16}+\frac{3\eta}{16}+\frac{\eta}{16}
\le 3\eta,
\]
where \(2\eta\) is the explicit projection/binomial error and the three
fractions come respectively from \(r\), the corrected \(q-q^2\) term, and the
linear \(K\)-term in (6).

It remains to compare \(C_0\) with the advertised negative constant.  Let
\(d=x_R-x_B\ge0\).  Since \(t=x_R+x_B-1\), we have
\[
x_B=\frac{1+t-d}{2}.
\]
Substituting the values of \(\alpha\) and \(q_0\) gives
\[
\begin{aligned}
C_0
&=\frac t2-\frac{1+\varepsilon}{2}t
-\frac{1+\varepsilon}{2}\cdot\frac{1}{2(1+\varepsilon)}(1-x_B)
+\frac{1+\varepsilon}{2}\cdot\frac{1}{4(1+\varepsilon)^2}
+\frac{1+\varepsilon}{2}\varepsilon^3\\
&=-\frac{\varepsilon t}{2}-\frac{1-x_B}{4}
+\frac{1}{8(1+\varepsilon)}
+\frac{(1+\varepsilon)\varepsilon^3}{2}.
\end{aligned}
\]
Now substitute \(x_B=(1+t-d)/2\).  This gives the exact expression
\[
C_0
=-\frac{\varepsilon}{8(1+\varepsilon)}
-\frac d8+\left(\frac18-\frac{\varepsilon}{2}\right)t
+\frac{(1+\varepsilon)\varepsilon^3}{2}. \tag{10}
\]
The middle-regime hypothesis gives \(t<\varepsilon/2\), and
\(0<\varepsilon<1/16\) makes \(1/8-\varepsilon/2>0\).  Therefore, using
\(d\ge0\),
\[
C_0\le
-\frac{\varepsilon}{8(1+\varepsilon)}
+\frac{\varepsilon}{16}
+\frac{(1+\varepsilon)\varepsilon^3}{2}
\le
\frac{\varepsilon(-1+\varepsilon+22\varepsilon^2)}
     {16(1+\varepsilon)}. \tag{11}
\]
For the last inequality, move the left side to the right side.  The difference
is
\[
\frac{\varepsilon^3(7-8\varepsilon-4\varepsilon^2)}
     {8(1+\varepsilon)},
\]
because, over the common denominator \(16(1+\varepsilon)\),
\[
\begin{aligned}
&\varepsilon(-1+\varepsilon+22\varepsilon^2)
-\left(-2\varepsilon+\varepsilon(1+\varepsilon)
  +8(1+\varepsilon)^2\varepsilon^3\right)\\
&\qquad =
14\varepsilon^3-16\varepsilon^4-8\varepsilon^5
=2\varepsilon^3(7-8\varepsilon-4\varepsilon^2).
\end{aligned}
\]
The denominator is positive since \(1+\varepsilon>0\).  Also
\(\varepsilon^3>0\), and the hierarchy bound \(\varepsilon<1/16\) gives
\[
8\varepsilon<\frac12,\qquad
4\varepsilon^2<\frac1{64},
\]
so
\[
7-8\varepsilon-4\varepsilon^2
>
7-\frac12-\frac1{64}
=\frac{415}{64}>0.
\]
Thus the displayed difference is nonnegative.  Combining (11) with (9)
proves the desired inequality.
\end{proof}
